% \chapter{Abstract} \label{abstract}
\begin{abstract}
\textsc{Sparse-Matrix Vector Multiplication} (SpMV) is a computational kernel which arise in many applications of scientific computing. It is a kernel that has an inherently low computational intensity, and its performance is therefore bound by the memory bandwidth of the architecture it is run on. The performance of SpMV can be improved through the use of both shared and distributed memory parallelization. Parallellizing SpMV in a distributed memory environment requires careful attention to managing the communication patterns which arise when a matrix is distributed on multiple nodes, in order to reduce the memory traffic across the interconnects.
\medskip

This thesis aims to investigate the impact that five increasingly selective communication strategies have on the performance of SpMV. To achieve this, these communication strategies have been tested on a set of parallel hardware that include both dual-socketed and single-socketed compute nodes. The configurations used on these compute nodes include both single-node non-shared memory, and a plethora of multi-node shared and distributed memory configurations. 
\medskip

The findings indicate from the experiments performed in this thesis indicate that....


% \textsc{Sparse Matrix Vector Multiplication} is an important kernel used in scientific computing. It is a problem that lends itself well to parlallization. The problem is bounded by, and scales with the memory bandwidth of the system. Therefore in order to efficiently perform \textit{SpMV} on large distributed memory systems, it is important to reduce the communication between nodes, in order to extract as much as possible out of the memory bandwidth of the system.

% This thesis aims to investigate the results of \textit{SpMV} when ran using different communication strategies.
\end{abstract}

